\section{Tag 4 — Bündelfehler, Interleaving und CRC}

\subsection*{Bleistiftübung 1 — Hamming bei Bündelfehlern}

(a) Bei zwei benachbarten Fehlern wird der $(7,4)$-Hamming-Code immer
\textbf{falsch korrigieren} — er denkt, es war ein Einzelfehler an
einer dritten, falschen Position, und kippt dort ein Bit. Resultat:
3 Bits falsch, Datenwort fast immer kaputt.

(b) Bei 3 Bündelfehlern: ähnlich, der Decoder produziert Müll. Mit
jeder weiteren Fehlerzahl wird das Ergebnis fast unkorreliert mit dem
Original.

(c) Hintereinanderschalten von Hamming-Codes löst das Problem
\textbf{nicht}. Schauen wir uns das genauer an: 4 Codewörter à 7 Bit
ergeben 28 Bits. Ein Bündelfehler von 3 Bits zwischen Wort 2 und Wort
3 trifft zum Beispiel die letzten 2 Bits von Wort 2 und das erste Bit
von Wort 3. Wort 2 hat also 2 Fehler → unrettbar.
Hintereinanderschalten reicht nicht; wir brauchen \textbf{Interleaving}.

\subsection*{Bleistiftübung 2 — SECDED}

(a) Codewort \texttt{0 1 1 0 0 1 1} hat 4 Einsen → Anzahl ist gerade →
$P = \mathbf{0}$. Vollständiges $(8,4)$-Codewort: \textbf{\texttt{0 1
1 0 0 1 1 0}}.

(b) Empfangen: \texttt{0 1 0 0 0 1 1 0} (Bit 3 ist gekippt verglichen
mit dem Original).

Hamming-Prüfsummen:
\begin{itemize}
\item $s_1 = b_1+b_3+b_5+b_7 \bmod 2 = 0+0+0+1 = 1$
\item $s_2 = b_2+b_3+b_6+b_7 \bmod 2 = 1+0+1+1 = 3 \bmod 2 = 1$
\item $s_3 = b_4+b_5+b_6+b_7 \bmod 2 = 0+0+1+1 = 2 \bmod 2 = 0$
\end{itemize}
$s_3\, s_2\, s_1 = 011$ binär $= 3$ → Hamming meldet Fehler an Position
3.

Gesamt-Parität des Empfangenen: $0+1+0+0+0+1+1+0 = 3$ (ungerade) →
Gesamt-Parität ist falsch.

Hamming-Prüfung \textbf{Fehler erkannt} + Gesamt-Parität \textbf{falsch}
→ \textbf{1 Fehler, korrigierbar} an Position 3. \checkmark

(c) Empfangen: \texttt{1 1 0 0 0 0 1 0}. Verglichen mit \texttt{0 1 1
0 0 1 1 0} sind Bits 1, 3, 6 gekippt — 3 Fehler. Aber wir tun so, als
wüssten wir das nicht.

Hamming-Prüfsummen:
\begin{itemize}
\item $s_1 = 1+0+0+1 = 2 \bmod 2 = 0$
\item $s_2 = 1+0+0+1 = 2 \bmod 2 = 0$
\item $s_3 = 0+0+0+1 = 1$
\end{itemize}
$s_3\, s_2\, s_1 = 100$ binär $= 4$ → Hamming meldet Fehler an Position
4.

Gesamt-Parität: $1+1+0+0+0+0+1+0 = 3$ (ungerade) → falsch.

Hamming-Prüfung \textbf{Fehler erkannt} + Gesamt-Parität \textbf{falsch}
→ SECDED denkt: 1 Fehler an Position 4. \textbf{Aber das ist falsch!}
Tatsächlich waren es 3 Fehler. SECDED korrigiert hier zur falschen
Antwort.

\begin{quote}
SECDED garantiert: 1 Fehler korrigiert, 2 Fehler erkannt, 3 oder mehr
Fehler — keine Garantie. Genau wie Hamming-Code: jenseits der
Garantiegrenze wird's gefährlich.
\end{quote}

\subsection*{Bleistiftübung 3 — Interleaving auf Karopapier}

(a) Hintereinander: \texttt{1010 0111 1100 0011}. Das sind 16 Bits.

(b) Tabelle und spaltenweise lesen:

\begin{verbatim}
    Sp1 Sp2 Sp3 Sp4
A:   1   0   1   0
B:   0   1   1   1
C:   1   1   0   0
D:   0   0   1   1
\end{verbatim}

Spaltenweise: \texttt{1010 0110 1101 0101}. (Spalte 1 = $1,0,1,0$,
Spalte 2 = $0,1,1,0$, …)

(c) Bündelfehler kippt Bits an Position 5–8: das ist Spalte 2, also je
das zweite Bit jedes Codeworts.

Original Spalte 2: $0,1,1,0$. Gestört: $1,0,0,1$.

Zurück in Tabelle:

\begin{verbatim}
A:   1   1   1   0   (vorher 1010)   → 1 Fehler in Position 2
B:   0   0   1   1   (vorher 0111)   → 1 Fehler in Position 2
C:   1   0   0   0   (vorher 1100)   → 1 Fehler in Position 2
D:   0   1   1   1   (vorher 0011)   → 1 Fehler in Position 2
\end{verbatim}

Jedes Codewort hat genau \textbf{1 Fehler}. Wenn die Codewörter
Hamming-codiert wären, könnte jeder Decoder seinen eigenen
Einzelfehler korrigieren!

(d) Ohne Interleaving würde der 4-Bit-Bündelfehler an Position 5–8
das gesamte Codewort B zerstören (4 Fehler in B, 0 in A/C/D). Hamming
kann B nicht retten.

\subsection*{Aufgabe 1.1 — Interleaving-Test}

Erwartete Ausgabe: alle 4 Datenwörter werden korrekt rekonstruiert.
Wenn du den Bündelfehler auf 5 Bits oder mehr erweiterst, könnte ein
Codewort zwei Fehler bekommen → dann scheitert die Korrektur. Ab 8
Bündelfehler wird es definitiv kritisch.

\subsection*{Bleistiftübung 4 — CRC}

(a) \texttt{10110000} (5 Datenbits + 3 angehängte Nullen).

(b) Polynomdivision $\bmod 2$:

\begin{verbatim}
   1 0 1 1 0 0 0 0
   1 0 1 1                ← G·1
   ─────────
   0 0 0 0 0 0 0 0
         (führende 0en überspringen, weiter unten)

   1 0 1 1 0 0 0 0
   1 0 1 1
   ───────
       0 1 1 0 0 0
       (führendes 0, eine Stelle weiter)

       1 1 0 0 0
       1 0 1 1                ← G·1 (an passender Stelle)
       ───────
       0 1 1 1 0
         1 1 1 0
         1 0 1 1               ← G·1
         ───────
           1 0 1
\end{verbatim}

Rest: \texttt{1 0 1}.

(c) CRC = \texttt{1 0 1}. Vollständig gesendetes Wort:
\texttt{1 0 1 1 0 1 0 1}.

(d) Wenn man \texttt{10110101} durch \texttt{1011} teilt, sollte der
Rest \textbf{0} sein — das ist die Eigenschaft, die zur CRC-Prüfung
genutzt wird.

(Hinweis: Das ist auch der Trick, warum man auf der Senderseite die
\texttt{000} anhängt: damit der Rest beim Empfänger genau dann 0 ist,
wenn nichts gestört wurde.)

\subsection*{Aufgabe 2.1 / 2.2 — Brute-Force-Test}

Erwartete Ausgabe (sinngemäß):

\begin{verbatim}
248 erkannt, 7 nicht erkannt
Nicht erkannte Muster (Anzahl): 7
\end{verbatim}

Die 7 nicht erkannten Muster sind genau die Vielfachen von $G(x) =
\texttt{1011}$ selbst (außer dem Null-Muster, das ja gar kein Fehler
ist):

\begin{verbatim}
00001011, 00010110, 00101100, 01011000, 10110000,
00010101, ... (etc.)
\end{verbatim}

Allgemein gilt: ein CRC mit $r$ Prüfbits hat genau $2^{n-r}$ Vielfache
von $G(x)$ im $n$-Bit-Raum, also $2^{n-r} - 1$ nicht-triviale
unentdeckte Fehlermuster.

Bei guter Wahl von $G(x)$: keine 1-Bit-Fehler unentdeckt, keine
2-Bit-Fehler bis zur Maximallänge unentdeckt — erst ab 3 Fehlerbits
oder sehr großen Lücken werden überhaupt einige Muster nicht erkannt.

\subsection*{Antworten zu den drei Reflexionsfragen}

\begin{enumerate}
\item \textbf{Korrektur vs.\ Erkennung:} bei stark gestörten Kanälen
  ist Korrektur essenziell, weil Wiederholung teuer oder unmöglich ist
  (denk an Voyager — ein Re-Send dauert Stunden). In LANs/Internet
  ist Erkennung + Re-Send fast immer billiger. Die Faustregel: je
  teurer ein Re-Send, desto mehr Korrektur. Mobilfunk ist ein
  Mischmasch (Forward Error Correction für Schnelligkeit, plus
  Re-Send-Mechanismen oben drauf).

\item \textbf{Modulo für Bytes:} genau das bereiten wir morgen vor.
  Mit Bytes (0–255) braucht man eine andere Modulo-Arithmetik.
  Naheliegend wäre Modulo 256 — aber das hat schlechte mathematische
  Eigenschaften (256 ist nicht prim und 2 hat dort keinen Inversen).
  Stattdessen nutzt man \textbf{Polynomarithmetik modulo eines
  irreduziblen Polynoms} über GF(2). Klingt furchterregend, ist aber
  konzeptionell genau wie ISBN-10 modulo 11 — nur eben mit Polynomen.

\item \textbf{Mehr/weniger Information:} CRC braucht weniger
  Information, weil es nur „ja/nein“-Entscheidungen treffen muss.
  Korrektur braucht mehr, weil es zusätzlich die Position des Fehlers
  identifizieren muss. Daraus folgt: Erkennung kann mit weniger
  Redundanz erreicht werden als Korrektur — wenn man Re-Sends
  akzeptieren kann.
\end{enumerate}
